% Generated draft: requires mathematical and source review.
\documentclass[11pt]{article}
\usepackage[margin=25mm]{geometry}
\usepackage{fontspec}
\setmainfont{Latin Modern Roman}
\usepackage{amsmath,amssymb,mathtools,unicode-math}
\setmathfont{Latin Modern Math}
\usepackage{enumitem,needspace,xurl,hyperref}
\setlength{\parindent}{0pt}
\setlength{\parskip}{5pt}
\setlength{\emergencystretch}{3em}
\newcommand{\N}{\mathbb{N}}
\newcommand{\Z}{\mathbb{Z}}
\newcommand{\Q}{\mathbb{Q}}
\newcommand{\R}{\mathbb{R}}
\newcommand{\C}{\mathbb{C}}
\newcommand{\F}{\mathbb{F}}
\newcommand{\A}{\mathbb{A}}
\newcommand{\PP}{\mathbb{P}}
\newcommand{\E}{\mathbb{E}}
\newcommand{\eps}{\varepsilon}
\newcommand{\dd}{\,\mathrm{d}}
\begin{document}
\section*{UnsolvedMath ID: 801161}
\textbf{Source problem code:} catalogue entry
\subsection{Definitions and mathematical statement}
A Poisson bracket on a commutative algebra \(\mathcal A\) is a bilinear, skew-symmetric operation \(\{\cdot,\cdot\}\) satisfying the Leibniz rule and the Jacobi identity. Two Poisson structures \(P,Q\) on the same underlying space are compatible if every scalar linear combination \(aP+bQ\) is again a Poisson structure, where the scalars belong to the scalar field of the chosen setting.

The supplied record does not define the class of \emph{geometric Poisson structures}, its notion of differential order, or the subclass called \emph{Darboux--Poisson}. These terms are retained as source terminology; the algebra, scalar field, underlying geometric space, and admissible regularity are not specified here.

\textbf{Problem (Bolsinov, Konyaev, Matveev, Problem 14).} For a geometric Poisson structure \(P\) of third order that is not Darboux--Poisson, describe all geometric Poisson structures \(Q\) of first order compatible with \(P\). In notation, once the source's geometric setting and terminology have been fixed, describe
\[
\mathcal C_1(P)=\bigl\{Q: Q\text{ is a first-order geometric Poisson structure compatible with }P\bigr\}.
\]
Determine also how large a compatible pencil of such first-order structures can be.

For a family \((Q_i)_{i\in I}\) of such structures, mutual compatibility means that every finite scalar linear combination of its members is Poisson. The dimension of its linear span is one possible precise measure of size. However, the supplied record does not specify whether its word ``pencil'' means a two-generator linear pencil or a more general mutually compatible linear family, or whether the requested bound is for fixed \(P\) or over all admissible \(P\). Compatibility of each \(Q_i\) with \(P\) alone does not impose mutual compatibility among the \(Q_i\). These qualifications must be resolved before the size question has a unique precise formulation. Nonzero linear combinations may also have lower order or vanish, so an exact-order requirement for every member of a pencil would need clarification.
\subsection{Short English statement}
For a third-order geometric Poisson structure that is not Darboux--Poisson, classify all compatible first-order geometric Poisson structures and determine the possible size of compatible pencils formed from them. The supplied excerpt leaves the geometric setting, Darboux--Poisson terminology, and intended meaning of pencil size unspecified.
\subsection{Sources}
\begin{itemize}
\item Local UnsolvedMath registry, supplied complete record with permanent ID \(801161\), source problem code \texttt{catalogue entry}, titled ``Nijenhuis Geometry and Integrable Systems --- Problem 14.'' The registry attributes the problem to Bolsinov, Konyaev, and Matveev, with proposed year 2024. Its status label is \texttt{open}; this is supplied source metadata, not an independently verified mathematical or literature assessment.
\item Bolsinov, Konyaev, Matveev, \emph{Research Problems on Nijenhuis Geometry and Integrable Systems} (2024), Problem 14. Supplied source URL: \url{https://arxiv.org/abs/2410.04276}. The record reports access on 2026-10-16 and extraction by \texttt{arxiv\_latex}; these are supplied provenance statements, not a claim of independent access or verification.
\end{itemize}
\end{document}
